% \chapter{Messdaten}
\setcounter{chapter}{2}
\addcontentsline{toc}{chapter}{\thechapter \, Messdaten} % Damit es im Inhaltsverzeichnis erscheint
\setcounter{table}{0}  % Setze den Tabellenzähler zurück
\label{ch:Protokoll}

\thispagestyle{empty}

\includepdf[
  pages=-,               
  pagecommand={\thispagestyle{empty}} 
]{Versuche/\versuchsnummer/Auswertung/Messprotokoll.pdf}



\addcontentsline{lot}{table}{\protect\numberline{\thechapter I} Entwicklung der Kreiselfrequenz}
\addcontentsline{lot}{table}{\protect\numberline{\thechapter II} Präzissionsdauer unter verschiedenen Auslenkungswinkeln}
\addcontentsline{lot}{table}{\protect\numberline{\thechapter III} Präzissionsdauer unter verschiedenen Massen und Massenposition}
\addcontentsline{lot}{table}{\protect\numberline{\thechapter IV} Messung von 10 Farbwechseln}
\addcontentsline{lot}{table}{\protect\numberline{\thechapter V} Zeitliche Entwicklung von Eigen- und Nutationsfrequenz}